%--------------------------------------------------------------------------------
%
%
%--------------------------------------------------------------------------------
\pagebreak
\section*{B branch inconditional}
\addcontentsline{toc}{section}{B branch unconditional}

\instructionentry{Syntax}{
  \texttt{B[.G] target} \\
}

\instructionentry{Format}{

	The B instruction uses the ... \\
	
  	\begin{flushleft}
    	\begin{tikzpicture}[scale=0.7, transform shape]
        
        % \draw[help lines, gray!70, dashed] (0,0) grid(16,1.5);
        
        \node[  tsRectangle, 
                minimum width=16cm,
                minimum height=1cm,
                text centered,
                fill=white!50] (blocks) at (8,1) { };
                
     	\draw[-] (1,0.5) -- (1, 1.5);
     	\draw[-] (3,0.5) -- (3, 1.5);
     	\draw[-] (5,0.5) -- (5, 1.5);
     	\draw[-] (6,0.5) -- (6, 1.5);
     	\draw[-] (6.5,0.5) -- (6.5, 1.5);
     	       	
     	\node at (0.5,0.25) {2};
    	\node at (2,0.25) {4};
    	\node at (4,0.25) {4};
    	\node at (5.75,0.25) {3};
   	 	\node at (11.5,0.25) {19};
   	 	
   	 	\node at (0.5,1) {\small 0};
		\node at (2,1) {\small \OpCodeCMP};
		\node at (4,1) {\small Reg R};
		\node at (5.5,1) {\small 0};
		\node at (6.25,1) {\small 1G};
		\node at (11.5,1) {\small target};
   	 	
       \end{tikzpicture}
	\end{flushleft}
}

\instructionentry{Description}{

The \texttt{B} instruction ...

	
}

\instructionentry{Operation}{

  	\texttt{RegR <- cmpOp( cond, RegB, RegA )} {(register format)}\\
  	\smallskip
  	\texttt{RegR <- cmpOp( cond, RegB, immOperand( ))} {(immediate format)}\\
   	\smallskip
  	\texttt{RegR <- cmpOp( cond, RegR, memOperand( ))} {(indexed formats)}\\
}

\instructionentry{Exceptions}{

	Data Alignment Trap\\
	Memory Access Violation Trap\\
	Memory Protection Violation Trap\\
	Data TLB Miss Trap
}

\instructionentry{Notes}{

	None. 
}




